% Based on templates/lecture-example.tex; source examples, not check questions.
% Source: Virtual Memory.pdf, slides 8.20--8.24, 8.27 and 8.32.
% All states and fault counts independently simulated; Clock uses new-page R=0.
\begingroup
\renewcommand{\recordid}[1]{}
\renewcommand{\relatedtopic}[2]{\hyperref[topic:#1]{#2}}
\exercisegroup{SC2005-L18-EX}{第 18 讲：页面置换算法的同序列对照}

\examplemeta
  {SC2005-L18-EX}
  {2026-10-22}
  {Virtual Memory 8.20--8.24、8.27、8.32；共用讲义访问序列}
  {页框状态与缺页次数已逐步核对}
  {已完成学习与笔记审阅}

\linespread{1.15}\selectfont
\setlength{\parskip}{0.4em}

\textbf{共同条件。}页面访问序列为
\[
  1,\ 2,\ 3,\ 4,\ 1,\ 2,\ 5,\ 1,\ 2,\ 3,\ 4,\ 5.
\]
初始页框均为空；首次装入也算缺页。下表 F 表示缺页，H 表示命中，$-$ 表示空页框；状态均为处理本次访问\emph{之后}的结果。除特别说明外，列表位置固定对应物理页框，不表示淘汰优先级。

\exercisequestion{SC2005-L18-E01}{FIFO：三个页框与四个页框}
\textbf{题目与来源。}8.20--8.21：分别使用三个和四个页框执行上述序列，计算 FIFO 缺页次数，判断是否存在 Belady 异常。对应\relatedtopic{SC2005-L18-T01}{FIFO 与 Belady 异常}。

\begin{checkedsolution}
FIFO 命中不更新装入顺序；缺页且已满时淘汰最早装入的驻留页。
{\small
\begin{longtable}{@{}rrclcl@{}}
\toprule
步 & 页 & 三页框 & 固定页框状态 & 四页框 & 固定页框状态 \\
\midrule
\endhead
\bottomrule
\endfoot
1  & 1 & F & $[1,-,-]$ & F & $[1,-,-,-]$ \\
2  & 2 & F & $[1,2,-]$ & F & $[1,2,-,-]$ \\
3  & 3 & F & $[1,2,3]$ & F & $[1,2,3,-]$ \\
4  & 4 & F & $[4,2,3]$ & F & $[1,2,3,4]$ \\
5  & 1 & F & $[4,1,3]$ & H & $[1,2,3,4]$ \\
6  & 2 & F & $[4,1,2]$ & H & $[1,2,3,4]$ \\
7  & 5 & F & $[5,1,2]$ & F & $[5,2,3,4]$ \\
8  & 1 & H & $[5,1,2]$ & F & $[5,1,3,4]$ \\
9  & 2 & H & $[5,1,2]$ & F & $[5,1,2,4]$ \\
10 & 3 & F & $[5,3,2]$ & F & $[5,1,2,3]$ \\
11 & 4 & F & $[5,3,4]$ & F & $[4,1,2,3]$ \\
12 & 5 & H & $[5,3,4]$ & F & $[4,5,2,3]$ \\
\end{longtable}}
因此三页框缺页 $9$ 次，四页框缺页 $10$ 次。第 $7$ 步后，三页框为 $\{5,1,2\}$，四页框为 $\{5,2,3,4\}$；前者并不是后者的子集。四页框方案随后访问 $1,2,3,4,5$ 连续缺页，展示了 FIFO 的 Belady 异常。
\end{checkedsolution}

\exercisequestion{SC2005-L18-E02}{四个页框下，OPT 与 LRU 如何选择}
\textbf{题目与来源。}8.22、8.24、8.27：使用四个初始为空的页框，分别执行 OPT 与 LRU，解释 victim 的选择并计算缺页次数。对应\relatedtopic{SC2005-L18-T02}{OPT}、\relatedtopic{SC2005-L18-T03}{LRU}。

\begin{checkedsolution}
两者前四步均缺页，形成 $[1,2,3,4]$；第 $5,6$ 步访问 $1,2$ 均命中。此时 LRU 从最久未访问到最近访问的顺序为 $[3,4,1,2]$，而 OPT 要查看尚未发生的引用。

\textbf{第 $7$ 步访问 $5$。}未来依次为 $1,2,3,4,5$，OPT 淘汰下一次使用最晚的 $4$；LRU 淘汰上一次访问最早的 $3$。后六步如下，圆括号内为淘汰的页号：
\begin{center}
\small
\begin{tabular}{@{}rrclcl@{}}
\toprule
步 & 页 & OPT & 固定页框状态 & LRU & 固定页框状态 \\
\midrule
7  & 5 & F$(4)$ & $[1,2,3,5]$ & F$(3)$ & $[1,2,5,4]$ \\
8  & 1 & H & $[1,2,3,5]$ & H & $[1,2,5,4]$ \\
9  & 2 & H & $[1,2,3,5]$ & H & $[1,2,5,4]$ \\
10 & 3 & H & $[1,2,3,5]$ & F$(4)$ & $[1,2,5,3]$ \\
11 & 4 & F$(1)$ & $[4,2,3,5]$ & F$(5)$ & $[1,2,4,3]$ \\
12 & 5 & H & $[4,2,3,5]$ & F$(1)$ & $[5,2,4,3]$ \\
\bottomrule
\end{tabular}
\end{center}
OPT 在第 $11$ 步面对的剩余引用只有 $5$；页 $1,2,3$ 均不再使用，淘汰任意一个都可以，表中沿用讲义选择 $1$。LRU 第 $10,11,12$ 步依次淘汰 $4,5,1$，依据每页最近一次访问的位置。因此
\[
 \boxed{\mathrm{OPT}:4+2=6\text{ 次}},\qquad
 \boxed{\mathrm{LRU}:4+4=8\text{ 次}}.
\]

\textbf{用栈复核 LRU（8.27）。}顶部在左，装满后为 $[4,3,2,1]$；命中 $1$ 后为 $[1,4,3,2]$，命中 $2$ 后为 $[2,1,4,3]$，装入 $5$ 时淘汰底部 $3$，得到 $[5,2,1,4]$。这一列表记录最近使用顺序，与上表固定物理页框的列法不同。
\end{checkedsolution}

\exercisequestion{SC2005-L18-E03}{Clock：同时追踪页面、引用位和指针}
\textbf{题目与来源。}8.32：四个初始为空的页框，扫描顺序 $F_1\to F_2\to F_3\to F_4\to F_1$。按原图，新装入页 $R=0$，随后命中置 $R=1$；替换后指针前进，命中时不动。计算缺页次数。对应\relatedtopic{SC2005-L18-T04}{Clock}。

\begin{checkedsolution}
前四步依次装满，缺页四次，所有 $R=0$，指针回到 $F_1$；第 $5,6$ 步命中页 $1,2$，只将其引用位置 $1$。下表每格为“页号$^{R}$”，“下一指针”是处理本步后的位置：
{\small
\begin{longtable}{@{}rrcccccc@{}}
\toprule
步 & 页 & $F_1$ & $F_2$ & $F_3$ & $F_4$ & 下一指针 & 结果 \\
\midrule
\endhead
\bottomrule
\endfoot
1  & 1 & $1^0$ & $-$ & $-$ & $-$ & $F_2$ & F \\
2  & 2 & $1^0$ & $2^0$ & $-$ & $-$ & $F_3$ & F \\
3  & 3 & $1^0$ & $2^0$ & $3^0$ & $-$ & $F_4$ & F \\
4  & 4 & $1^0$ & $2^0$ & $3^0$ & $4^0$ & $F_1$ & F \\
5  & 1 & $1^1$ & $2^0$ & $3^0$ & $4^0$ & $F_1$ & H \\
6  & 2 & $1^1$ & $2^1$ & $3^0$ & $4^0$ & $F_1$ & H \\
7  & 5 & $1^0$ & $2^0$ & $5^0$ & $4^0$ & $F_4$ & F \\
8  & 1 & $1^1$ & $2^0$ & $5^0$ & $4^0$ & $F_4$ & H \\
9  & 2 & $1^1$ & $2^1$ & $5^0$ & $4^0$ & $F_4$ & H \\
10 & 3 & $1^1$ & $2^1$ & $5^0$ & $3^0$ & $F_1$ & F \\
11 & 4 & $1^0$ & $2^0$ & $4^0$ & $3^0$ & $F_4$ & F \\
12 & 5 & $1^0$ & $2^0$ & $4^0$ & $5^0$ & $F_1$ & F \\
\end{longtable}}
第 $7$ 步先清零并跳过页 $1,2$，淘汰页 $3$；第 $10$ 步直接淘汰指针处的页 $4$；第 $11$ 步再次清零并跳过 $1,2$，淘汰页 $5$；最后淘汰页 $3$。缺页发生在步骤 $1,2,3,4,7,10,11,12$，共 $\boxed{8}$ 次。

LRU 在最后一步淘汰页 $1$，Clock 淘汰页 $3$：本例缺页次数相等，但选择不同，引用位并不能恢复精确的最近访问顺序。
\end{checkedsolution}
\endgroup
